\documentclass[11pt]{article}
\usepackage[utf8]{inputenc}
\usepackage[T1]{fontenc}
\usepackage[spanish]{babel}
\usepackage{amsmath, amssymb, bm}
\usepackage{geometry}
\geometry{margin=2.5cm}

\title{Ridge/LS 03: algebra lineal, rango, invertibilidad y diagonales}
\author{}
\date{}

\begin{document}
\maketitle

\section{Por que hace falta algebra lineal}

La expresion:
\[
    (A^H A + \lambda I)^{-1}A^H
\]
solo tiene sentido si entendemos algunas ideas de algebra lineal:

\begin{itemize}
    \item matrices complejas;
    \item transpuesta conjugada;
    \item rango;
    \item matrices Hermitianas;
    \item matrices definidas positivas;
    \item invertibilidad;
    \item eigenvalues;
    \item condicionamiento numerico.
\end{itemize}

\section{Matrices complejas y transpuesta conjugada}

En comunicaciones digitales aparecen exponenciales complejas:
\[
    e^{-j2\pi k\ell/N}.
\]

Por eso \(A\) es compleja.

Para matrices complejas, la transpuesta adecuada es:
\[
    A^H = \overline{A}^T.
\]

Primero se conjuga y luego se transpone.

Ejemplo:
\[
A =
\begin{bmatrix}
    1+j & 2\\
    -j & 3-j
\end{bmatrix}
\]

entonces:
\[
A^H =
\begin{bmatrix}
    1-j & j\\
    2 & 3+j
\end{bmatrix}.
\]

\section{Rango}

El rango de \(A\) indica cuantas columnas realmente independientes tiene.

Si:
\[
    A \in \mathbb{C}^{M\times L},
\]
entonces el maximo rango posible es:
\[
    \min(M,L).
\]

Para estimar \(L\) taps con LS, se necesita:
\[
    \mathrm{rank}(A)=L.
\]

Esto implica que debe haber al menos tantas observaciones utiles como taps:
\[
    M \ge L.
\]

En el codigo aparece una proteccion:
\[
    \texttt{tap\_count} \le \texttt{len(known\_indices)}.
\]

Eso evita intentar estimar mas taps que ecuaciones disponibles.

\section{La matriz \(A^H A\)}

Definimos:
\[
    G = A^H A.
\]

Si \(A\) tiene dimensiones \(M\times L\), entonces:
\[
    G \in \mathbb{C}^{L\times L}.
\]

Es decir, \(G\) es cuadrada aunque \(A\) no lo sea.

\section{Diagonal de \(A^H A\)}

El elemento diagonal \(G_{ii}\) es:
\[
    G_{ii} = \bm a_i^H \bm a_i = \|\bm a_i\|_2^2.
\]

Aqui \(\bm a_i\) es la columna \(i\)-esima de \(A\).

Por tanto, la diagonal de \(A^H A\) contiene energias de columnas.

\section{Elementos fuera de la diagonal}

El elemento \(G_{ij}\) es:
\[
    G_{ij} = \bm a_i^H \bm a_j.
\]

Esto es el producto interno entre dos columnas.

Si \(G_{ij}=0\), las columnas son ortogonales. Si \(|G_{ij}|\) es grande, las
columnas se parecen.

Cuando muchas columnas se parecen, \(A^H A\) se vuelve dificil de invertir.

\section{Matriz Hermitiana}

Una matriz \(G\) es Hermitiana si:
\[
    G^H = G.
\]

La matriz \(A^H A\) siempre es Hermitiana:
\[
    (A^H A)^H = A^H (A^H)^H = A^H A.
\]

Las matrices Hermitianas tienen eigenvalues reales. Esto es muy importante
para analizar estabilidad.

\section{Semidefinida positiva}

Para cualquier vector \(\bm x\):
\[
\begin{aligned}
    \bm x^H A^H A \bm x
    &=
    (A\bm x)^H(A\bm x)\\
    &=
    \|A\bm x\|_2^2\\
    &\ge 0.
\end{aligned}
\]

Por eso \(A^H A\) es semidefinida positiva.

Sus eigenvalues cumplen:
\[
    \mu_i \ge 0.
\]

\section{Definida positiva e invertibilidad}

Si \(A\) tiene rango columna completo:
\[
    \bm x \ne 0
    \quad \Rightarrow \quad
    A\bm x \ne 0.
\]

Entonces:
\[
    \bm x^H A^H A\bm x = \|A\bm x\|^2 > 0.
\]

En ese caso \(A^H A\) es definida positiva y todos sus eigenvalues son:
\[
    \mu_i > 0.
\]

Por tanto es invertible.

\section{Que hace \(\lambda I\) en la diagonal}

Si:
\[
G =
\begin{bmatrix}
g_{11} & g_{12}\\
g_{21} & g_{22}
\end{bmatrix},
\]
entonces:
\[
G+\lambda I =
\begin{bmatrix}
g_{11}+\lambda & g_{12}\\
g_{21} & g_{22}+\lambda
\end{bmatrix}.
\]

La regularizacion aumenta solo la diagonal.

Intuitivamente, hace que la matriz se parezca un poco mas a una matriz diagonal
dominante, lo cual suele mejorar la inversion numerica.

\section{Eigenvalues}

Si \(G=A^H A\) tiene eigenvalues:
\[
    \mu_1,\mu_2,\ldots,\mu_L,
\]
entonces:
\[
    G+\lambda I
\]
tiene eigenvalues:
\[
    \mu_1+\lambda,\mu_2+\lambda,\ldots,\mu_L+\lambda.
\]

Esto es crucial.

Si algun \(\mu_i\) era cero o muy pequeno, al sumar \(\lambda\) queda:
\[
    \mu_i+\lambda \ge \lambda.
\]

Por tanto, ya no hay eigenvalues menores que \(\lambda\).

\section{Condicionamiento}

El numero de condicion de una matriz definida positiva es:
\[
    \kappa(G) = \frac{\mu_{\max}}{\mu_{\min}}.
\]

Si \(\mu_{\min}\) es muy pequeno, \(\kappa(G)\) es muy grande y la inversion
amplifica errores.

Con regularizacion:
\[
    \kappa(G+\lambda I)
    =
    \frac{\mu_{\max}+\lambda}{\mu_{\min}+\lambda}.
\]

Esto suele ser menor que \(\kappa(G)\), especialmente si \(\mu_{\min}\) era muy
pequeno.

\section{Conclusion}

Desde el punto de vista de algebra lineal:
\[
    A^H A
\]
contiene la geometria de las columnas de \(A\), y:
\[
    \lambda I
\]
mejora la invertibilidad al aumentar todos los eigenvalues.

Por eso:
\[
    (A^H A + \lambda I)^{-1}A^H
\]
es mas estable que:
\[
    (A^H A)^{-1}A^H.
\]

\end{document}

